\documentclass[level=3,teacher]{kidscode}
%% Terminal and text-file listings: no Python keyword coloring.
\lstdefinestyle{kcshell}{basicstyle=\ttfamily\small\setstretch{1.0},language={},
  columns=fullflexible,keepspaces=true,breaklines=true,numbers=none,upquote=true,
  showstringspaces=false}
\begin{document}
\kctitle{Project Teacher Guide}{Build Something Big: milestones, rubric, reference solution and showcase}

\begin{tip}[In one minute]
Students build one larger project in weeks 9--12, alone or in pairs, choosing a grid game, a
generative art app or a data/AI app. Your job is to keep projects \textbf{small enough to finish},
check the milestones, unblock with questions rather than code, and run a warm showcase. The
reference solution, \textbf{Maze Muncher}, shows what a finished grid game looks like; keep its
code to yourself until the showcase, but do play it with the group at the kickoff.
\end{tip}

\setcounter{tocdepth}{1}
\tableofcontents

%% =====================================================================
\section{1. Overview and preparation}

\renewcommand{\arraystretch}{1.3}
\begin{longtable}{@{}>{\bfseries}p{1.6cm}>{\raggedright\arraybackslash}p{4.2cm}p{\dimexpr\linewidth-5.8cm-4\tabcolsep}@{}}
\toprule
\rowcolor{kcaccent!15} When & Milestone & In the session\\ \midrule\endhead
\bottomrule\endfoot
Week 9 & Kickoff $\rightarrow$ \textbf{M1 Pitch} (due end of week 9) & Last 13 minutes of the
week 9 lesson: play Maze Muncher together, hand out the brief, read the must-haves, choose
options. Pitch page is homework.\\ \midrule
Week 10 & \textbf{M2 Design + core classes} & CRC cards and class diagram on paper (about 11 min),
copy the starter, \texttt{git init}, class skeletons and first commits. Partner code review at the
end. Working methods and a first test are homework.\\ \midrule
Week 11 & \textbf{M3 Features + 5 tests + playtest}, then \textbf{M4 README + polish} & First half:
tests, 8-minute playtest swap and triage. Second half: fix the top playtest issue, write the
README, final commits.\\ \midrule
Week 12 & \textbf{Showcase} & Gallery walk and 2-minute talks, then the exit assessment and
certificates (see section 7).\\
\end{longtable}

\begin{materials}
  \item Print: project brief (one per student, in week 9); 2 playtest cards each (in the brief);
        CRC cards can be index cards or the brief's page B.
  \item On each Mac: a folder \texttt{Documents/coding-projects} (students copy
        \texttt{project/starter} into it and rename it in week 10), git set up with a nickname
        (Teacher Guide step 5.2).
  \item On your Mac only: \texttt{project/reference}. Run \texttt{maze\_muncher.py} once in Thonny
        so you know the controls (arrow keys; SPACE for the next level; R to restart).
  \item Sticky notes, a timer, and a \enquote{parking lot} sheet for nice-to-have ideas.
\end{materials}

%% =====================================================================
\section{2. Milestone walkthrough}

\subsection{M1 Pitch (end of week 9)}
\begin{teacherscript}
\enquote{For the next four weeks you are the boss of your own project. Real developers start
small and make it bigger step by step. Your pitch should be something you can finish in two
sessions, plus ideas for later in the parking lot.}
\end{teacherscript}
\textbf{Check before you say \enquote{go}:} the pitch names 3 must-have features (not 10); it says
which algorithm and data structure; it saves a file; the student can say the idea in one sentence;
it names a path.

\begin{teachernote}[Core path or Full path]
Two 60-minute studio sessions plus homework is not much time for a 9--11-year-old. The brief
(section 2, \enquote{Pick your path}) therefore offers a \textbf{Core path}: keep a template's
classes and add exactly one class, reuse an algorithm from class (flood fill from week 7,
recursion from the week 7 fractal tree, or nearest neighbor from week 9), add a best-score file
with the recipe in the brief, and count the template's example tests towards the five. Every
must-have is still met, and the rubric is the same: a Core project can reach \enquote{Got it}
(3 points) in every row, about 18 of 24. \enquote{Expert} needs the student's own ideas on top.
Recommend the Core path to most students, and to every student whose pitch is bigger than one
session's work. Concretely: (a) \texttt{Game} has a \texttt{Grid} and a \texttt{Hero}, plus a
recursive \texttt{count\_reachable()} flood fill; (b) \texttt{Tree(Shape)} with a recursive
\texttt{branch()}, added to \texttt{SHAPE\_CLASSES} so the gallery saves and loads it;
(c) \texttt{Guesser} has a \texttt{DataSet} and a \texttt{guess()} method that copies the genre of
the row with the smallest \texttt{distance()}. All three were built against the templates and pass
their tests.
\end{teachernote}
\textbf{Good question:} \enquote{If you only had one session, what would you build?} That answer
is feature 1. In week 10 each student copies the starter folder into
\texttt{Documents/coding-projects}, renames it and makes the first commit (the starter README
lists the exact Terminal lines).

\subsection{M2 Design + core classes (week 10)}
\begin{itemize}[itemsep=2pt]
  \item \textbf{Nouns become classes, verbs become methods.} Ask students to underline nouns in
        their pitch (dragon, grid, treasure) and circle verbs (move, collect, save).
  \item \textbf{CRC cards} (10 min): one card per class. If a card has more than about 5
        responsibilities, split the class. If a class has no responsibilities, drop it.
  \item \textbf{Inheritance or composition?} Ask \enquote{Is a Dragon \emph{a kind of}
        Character?} (inheritance) and \enquote{Does a Game \emph{have} a Grid?} (composition).
        Either one meets the requirement.
  \item \textbf{Build the core:} in the session, class skeletons (\texttt{\_\_init\_\_} and
        empty methods) with a commit after each class; at home, one or two methods that really
        work, a first test, and 3 more commits (homework 10).
\end{itemize}
\textbf{Check (end of the session):} at least 2 CRC cards with one is-a or has-a link, class
skeletons that run without errors, a repository with at least 2 commits.

\subsection{M3 Features + 5 tests + playtest (week 11, first half)}
\begin{itemize}[itemsep=2pt]
  \item Features are added one at a time: write or update a test, make it pass, commit.
  \item \textbf{Playtest swap} (8 min): each student plays someone else's project while the owner
        writes feedback cards (workbook Exercise 11.2; the brief's playtest card works too). With one student, you are the playtester; with an odd number, rotate.
  \item Each student circles \textbf{one} wish on their card to fix first.
\end{itemize}
\textbf{Check:} must-have features work; \texttt{python -m pytest -q} (in Thonny's system shell) shows 5 or more passing.

\subsection{M4 README + polish (week 11, second half)}
\begin{itemize}[itemsep=2pt]
  \item Fix the circled playtest wish. Run the tests again. Commit.
  \item Copy the README template from the starter \texttt{README.md}, fill in every section (15
        minutes is enough). The \enquote{design decision} section is rehearsal for the showcase.
  \item \texttt{git log --oneline}: count the commits (6+ now, 8+ by the showcase). Fewer? Usually the work is there
        but was committed in big lumps; from now on commit after each small change.
\end{itemize}

%% =====================================================================
\section{3. Unblocking tips}

Coach with a question first; give the hint only if the question does not help.

\begin{longtable}{@{}>{\raggedright\arraybackslash}p{3.6cm}>{\raggedright\arraybackslash}p{4.6cm}>{\raggedright\arraybackslash}p{\dimexpr\linewidth-8.2cm-4\tabcolsep}@{}}
\toprule
\rowcolor{kcaccent!15} Symptom & Ask & Hint\\ \midrule\endhead
\bottomrule\endfoot
The idea is huge & \enquote{What is the smallest version that is still fun?} & Three features now; the
rest goes on the parking-lot sheet. Finished beats perfect.\\ \midrule
\enquote{I don't know which classes} & \enquote{What things are in your program? What do they do?} &
Nouns $\rightarrow$ classes, verbs $\rightarrow$ methods. Start from the template's classes.\\ \midrule
No algorithm yet & \enquote{Is there anything to find, order, or repeat inside itself?} &
(a) BFS chaser, flood fill to reveal rooms; (b) recursive fractal, sort shapes by size before
drawing; (c) sort for a top 3, nearest-neighbor search for the closest example.\\ \midrule
Tests are hard because of turtle & \enquote{Which part decides, and which part draws?} & Move rules into a
class or function with no turtle (like \texttt{muncher\_model.py}); test that part.\\ \midrule
\texttt{ModuleNotFoundError} in tests & \enquote{What is your file called now?} & The test's
\texttt{import} must match the file name, without \texttt{.py}, and the file name must not contain
spaces or dashes.\\ \midrule
Few commits & \enquote{When did something last work?} & Commit after every passing test. Messages
start with a verb.\\ \midrule
\enquote{I broke everything} & \enquote{What did you change since your last commit?} & \texttt{git diff}
shows it; \texttt{git restore file.py} goes back to the last commit. Celebrate the safety net.\\ \midrule
No saved file & \enquote{What would you want to keep when you close the program?} & High scores,
levels, saved art (JSON), a results file. Create it if it is missing.\\ \midrule
Pair disagreement & \enquote{Which class could each of you own?} & Split by class; swap driver every 8
minutes; decide with a coin when stuck --- and write the decision in the README.\\ \midrule
Finished early & \enquote{What would surprise your playtester?} & Extension path in the brief: a
second level from a file, test-first, a settings JSON, k-nearest neighbors.\\
\end{longtable}

%% =====================================================================
\section{4. Rubric and scoring}

First check the \textbf{must-haves} (brief, section 2). Then score each criterion from 1 to 4
(Just starting = 1, Getting there = 2, Got it = 3, Expert = 4). Total out of 24 (six criteria
$\times$ 4 points; the lowest possible total is 6). Core-path and Full-path projects use the same
rubric.

\begingroup\small
\renewcommand{\arraystretch}{1.3}
\begin{longtable}{@{}>{\bfseries\raggedright\arraybackslash}p{2.2cm}>{\raggedright\arraybackslash}p{3.15cm}>{\raggedright\arraybackslash}p{3.15cm}>{\raggedright\arraybackslash}p{3.15cm}>{\raggedright\arraybackslash}p{\dimexpr\linewidth-11.65cm-8\tabcolsep}@{}}
\toprule
\rowcolor{kcaccent!15} Criterion & 1 Just starting & 2 Getting there & 3 Got it & 4 Expert\\ \midrule\endhead
\bottomrule\endfoot
Creativity & Copies an example with small changes & Changes an example noticeably & Own idea and
own style (art, theme, rules) & Something genuinely new or surprising; thoughtful for its audience\\ \midrule
Working code & Crashes during normal use & Runs; some features buggy or missing & All must-have
features work in a normal demo & Robust to unexpected input; no crashes in the playtest\\ \midrule
Design / OOP (L3.G1, G3) & One class, or classes that do not help & Two classes; responsibilities
mixed; data structure used thinly & Clear classes with one job each; inheritance or composition
used sensibly; suitable data structure & Design makes changes easy; student shows how (e.g.
adding a new enemy is one small class)\\ \midrule
Testing (L3.G8) & 0--2 tests & 3--4 tests, or some failing & 5+ passing tests of the logic & Edge
cases (empty, unreachable, missing file); a test caught a real bug\\ \midrule
Git history (L3.G4) & Under 4 commits & 4--7 commits or vague messages & 8+ commits, clear verb-first
messages & The log tells the story; small steps; tests committed with features\\ \midrule
Explaining design decisions (L3.G5) & Shows what it does & Explains how one part works & Explains
why a design choice was made & Compares with an alternative and explains the trade-off\\
\end{longtable}
\endgroup

\textbf{Scoring bands:} 21--24 Expert; 16--20 Got it; 11--15 Getting there; 6--10 Just
starting. Explaining is weighted the same as working code on purpose: a smaller project explained
with real understanding is a better Level 3 result than a big one the student cannot explain.

\textbf{How to check quickly:} in the student's project folder run \texttt{git log --oneline} (read
the messages, count the lines) and, in Thonny's system shell, \texttt{python -m pytest -q} (count
the passing tests); open
the README; ask \enquote{Why did you design it this way?}

%% =====================================================================
\section{5. Reference solution: Maze Muncher}

\subsection{What it is}
A grid maze game: the player (\textbf{P}, orange circle) eats every pellet (small black dot) while
a wandering ghost (\textbf{G}, sky-blue triangle) and a chasing ghost (\textbf{C}, vermillion
diamond) try to catch them. Each character has a \textbf{shape and a letter}, so nobody needs to
tell colors apart. Clearing level 1 unlocks level 2; the best five scores are saved. It meets
every must-have:

\begin{tabularx}{\linewidth}{@{}>{\bfseries}p{4.2cm}X@{}}
\toprule
Must-have & Where in Maze Muncher\\ \midrule
2+ classes, inheritance, composition & \texttt{Entity} $\rightarrow$ \texttt{Player}, \texttt{Ghost} $\rightarrow$ \texttt{Chaser}
(inheritance); \texttt{Game} has a \texttt{Maze}, a \texttt{Player} and ghosts (composition)\\
Data structure beyond a flat list & the maze is a 2D grid; BFS uses a queue (\texttt{deque}) and a dict\\
Named algorithm & breadth-first search in \texttt{Maze.shortest\_path}\\
Saved file & \texttt{high\_scores.txt} (\texttt{nickname,score} lines), created if missing\\
5+ tests & 16 model tests + 5 view tests\\
Accessible, private & shapes + letters; asks for a nickname, never a real name\\
\bottomrule
\end{tabularx}

\subsection{Files}
\begin{tabularx}{\linewidth}{@{}>{\ttfamily}p{4.6cm}X@{}}
muncher\_model.py & all the rules, no turtle: \texttt{Maze}, \texttt{Entity}, \texttt{Player},
\texttt{Ghost}, \texttt{Chaser}, \texttt{Game}, high-score functions\\
maze\_muncher.py & the turtle view: drawing, keys, the game clock (\texttt{MuncherApp})\\
levels/level1.txt & the first maze, as text (\texttt{level2.txt} is the second)\\
test\_muncher\_model.py & 16 tests of the rules\\
test\_maze\_muncher.py & 5 tests of the view on a fake screen\\
\end{tabularx}

\subsection{The design}
\begin{center}
\begin{tikzpicture}[
  cls/.style={draw, thick, rounded corners=2pt, align=left, font=\small, fill=kcsky!10, inner sep=4pt},
  isa/.style={-{Triangle[open,length=7pt,width=7pt]}, thick},
  hasa/.style={{Diamond[length=9pt,width=6pt]}-, thick}]
  \node[cls] (app) at (0,0) {\textbf{MuncherApp} (view)\\ screen, pens, game\\ redraw(), tick(), on\_key()};
  \node[cls] (game) at (6.4,0) {\textbf{Game}\\ maze, player, ghosts,\\ score, state\\ move\_player(), tick()};
  \node[cls] (maze) at (11.6,0) {\textbf{Maze}\\ grid, starts\\ is\_wall(), neighbors(),\\ shortest\_path()};
  \node[cls] (entity) at (6.4,-3.3) {\textbf{Entity}\\ pos, previous, symbol\\ try\_move()};
  \node[cls] (player) at (2.6,-6.0) {\textbf{Player}\\ symbol = "P"};
  \node[cls] (ghost) at (8.6,-6.0) {\textbf{Ghost}\\ symbol = "G"\\ choose\_step() (random)\\ take\_turn()};
  \node[cls] (chaser) at (8.6,-9.0) {\textbf{Chaser}\\ symbol = "C"\\ choose\_step() (BFS)};
  \draw[hasa] (app.east) -- node[above, font=\scriptsize] {has a} (game.west);
  \draw[hasa] (game.east) -- node[above, font=\scriptsize] {has a} (maze.west);
  \draw[hasa] (game.south) -- node[right, font=\scriptsize] {has} (entity.north);
  \draw[isa] (player.north) -- node[left, font=\scriptsize] {is a} (entity.south west);
  \draw[isa] (ghost.north) -- node[right, font=\scriptsize] {is a} (entity.south east);
  \draw[isa] (chaser.north) -- node[right, font=\scriptsize] {is a} (ghost.south);
\end{tikzpicture}
\end{center}

\textbf{Design decisions to point out at the showcase} (and to model the \enquote{explain your
design} criterion):
\begin{enumerate}[itemsep=2pt]
  \item \textbf{Logic and drawing are separate files.} \texttt{muncher\_model.py} never imports
        turtle, so all its rules are tested in a fraction of a second without opening a window.
  \item \textbf{Chaser inherits from Ghost and overrides only \texttt{choose\_step}.} The game clock
        calls \texttt{ghost.take\_turn(\ldots)} for every ghost and each kind decides its own way
        (polymorphism). Adding a new ghost type would be one small class.
  \item \textbf{Levels are text files.} Anyone can design a new maze in Thonny without touching
        code; a test checks that every pellet in every shipped level can be reached.
  \item \textbf{One game clock.} \texttt{MuncherApp} never starts a second \texttt{ontimer} loop
        (the \texttt{clock\_running} flag), a classic turtle-game bug.
\end{enumerate}

\subsection{The level file format}
One line per row. \texttt{\#} wall, \texttt{.} pellet, \texttt{P} player start (exactly one),
\texttt{G} wandering ghost, \texttt{C} chaser. Ghost squares also get a pellet underneath; short
rows are padded with walls.
\begin{tcolorbox}[kccode,left=6pt,title={\small\faFile*[regular]\ \texttt{levels/level1.txt}}]
\lstinputlisting[style=kcshell]{reference/levels/level1.txt}
\end{tcolorbox}

\subsection{Breadth-first search, step by step}
\texttt{Maze.shortest\_path(start, goal)} explores outwards in rings. Below, the number in each
square is how many steps it is from \textbf{P}; BFS reaches every square in order 1, 2, 3, 4.
When it first reaches the goal (here the square marked 4, 4 steps away), it walks back through
the \texttt{came\_from} dictionary to build the path (arrows).

\begin{center}
\begin{tikzpicture}[x=1cm,y=-1cm]
  \foreach \r/\c in {0/0,0/1,0/2,0/3,0/4,1/0,1/4,2/0,2/2,2/4,3/0,3/4,4/0,4/1,4/2,4/3,4/4}
    \draw[fill=kcblue, draw=black] (\c,\r) rectangle ++(1,1);
  \foreach \r/\c/\d in {1/1/P,1/2/1,1/3/2,2/1/1,2/3/3,3/1/2,3/2/3,3/3/4}
    \draw (\c,\r) rectangle ++(1,1) node[midway, font=\bfseries] {\d};
  \draw[very thick, -{Stealth[length=7pt]}] (1.5,1.75) -- (1.5,2.25);
  \draw[very thick, -{Stealth[length=7pt]}] (1.5,2.75) -- (1.5,3.25);
  \draw[very thick, -{Stealth[length=7pt]}] (1.75,3.5) -- (2.25,3.5);
  \draw[very thick, -{Stealth[length=7pt]}] (2.75,3.5) -- (3.25,3.5);
  \node[anchor=west, font=\small, align=left] at (5.4,2.5)
    {filled blue squares = walls (\texttt{\#})\\ numbers = steps from P\\ arrows = the path BFS returns\\
     goal = the square marked 4};
\end{tikzpicture}
\end{center}

The queue matters: taking squares from the \emph{front} (\texttt{popleft}) means all 1-step squares
are explored before any 2-step square. If you took them from the back (a stack), you would get
\emph{a} path, but not always the shortest one. The \texttt{Chaser} simply takes \texttt{path[1]},
the first step of that path, every tick.

\subsection{Tests}
\texttt{test\_muncher\_model.py} checks level parsing (size, starts, padding, exactly one P, every
pellet reachable), walls blocking moves, eating pellets once, BFS on a corridor, around a wall and
to an unreachable square, the chaser getting one step closer, the wandering ghost never walking
into walls or turning back in a corridor, winning, losing, and the high-score file (created if
missing, round trip, best five kept, broken lines skipped) using pytest's \texttt{tmp\_path}.
\texttt{test\_maze\_muncher.py} loads the view on a fake screen (\texttt{kidtest.load\_module}),
presses keys, runs the game clock and checks walls, letters, scores and saved high scores.

\subsection{The code}
\pyfile[basicstyle=\ttfamily\footnotesize\setstretch{1.0}]{reference/muncher_model.py}
\pyfile[basicstyle=\ttfamily\footnotesize\setstretch{1.0}]{reference/maze_muncher.py}
\pyfile[basicstyle=\ttfamily\footnotesize\setstretch{1.0}]{reference/test_muncher_model.py}
\pyfile[basicstyle=\ttfamily\footnotesize\setstretch{1.0}]{reference/test_maze_muncher.py}

%% =====================================================================
\section{6. What a good git history looks like}

This is the kind of \texttt{git log --oneline} you are hoping for (newest first). Every message
starts with a verb, each commit is one small step, and tests arrive together with the features
they check. Show it to students in week 10 as a model.

\needspace{20\baselineskip}
\begin{shellcode}
$ git log --oneline
6636dc1 Write README with controls and design decisions
aa8bb0f Add level2.txt and SPACE for next level
1a48e61 Draw letters P, G, C on characters for accessibility
19fe9fa Save top-5 high scores to high_scores.txt (+2 tests)
8871c26 Add win and lose states
d107ee0 Test BFS on small mazes, including unreachable goal
aef335f Add Chaser that uses BFS to find the player
b6ac07e Add Ghost that wanders without turning back
f48dbaf Eat pellets and keep score (+2 tests)
e79159e Stop the player walking through walls (+2 tests)
8ddc8c2 Add Player class and arrow-key movement
e02d04b Draw walls and pellets with turtle
1b438bc Add Maze class that reads levels/level1.txt
61ad517 Start Maze Muncher from the grid game starter
\end{shellcode}

\begin{commonmistakes}
  \item One giant commit at the end (\texttt{"finished project"}): the history tells no story and
        nothing can be undone in small steps.
  \item Messages that say nothing: \texttt{"stuff"}, \texttt{"fix"}, \texttt{"asdfg"}.
  \item Committing personal information or a real name. The nickname setup prevents the author
        line problem; check the files too.
\end{commonmistakes}

%% =====================================================================
\clearpage
\section{7. Showcase plan (week 12)}

\subsection{Running order (60 minutes)}
This is the order of the week 12 lesson plan, which has the full script: the showcase comes
first, while projects are fresh on screen, then the exit assessment, then the certificates. If
your group is calmer early on, you may give the exit assessment first (the lesson plan's
\enquote{short on time} note).

\begin{timeplan}
  \timeslot{0--4}{Set-up}{Each project open on a Mac, window large, Thonny font enlarged. Families
  sit where they can see. Hand out feedback cards.}
  \timeslot{4--12}{Model talk}{You show Maze Muncher and give the model 2-minute talk (Problem,
  Choice, Why, Next time), then show its \texttt{git log}.}
  \timeslot{12--34}{Showcase}{Gallery walk (6 min) with a star and a wish per project, then each
  student (or pair) gives a 2-minute talk using the workbook talk card, plus one question.}
  \timeslot{34--54}{Assessment}{Quiet, paper only. Collect the papers; score them later.}
  \timeslot{54--60}{Graduation}{Certificates (names written by hand), one specific compliment each,
  photo only if every family agrees.}
\end{timeplan}

\subsection{Design-decision talk and demo script}
\begin{teacherscript}[Prompt card for a nervous presenter]
\enquote{My project is called \underline{\hspace{2.5cm}}. It is for \underline{\hspace{2.5cm}}.
Let me show you \underline{\hspace{2.5cm}}. One decision I made was \underline{\hspace{2.5cm}}
instead of \underline{\hspace{2.5cm}}, because \underline{\hspace{2.5cm}}. A bug I fixed was
\underline{\hspace{2.5cm}}. Next I would add \underline{\hspace{2.5cm}}.}
\end{teacherscript}
Demo tips: start the program \emph{before} the talk; show the best feature first; if it crashes,
smile and say \enquote{and that is the bug I will fix next} --- that is real software development.

\subsection{Feedback cards}
Each audience member fills one card per presenter: \textbf{Star:} something great.
\textbf{Star:} something clever in the design. \textbf{Wish:} one idea for next time. Read the
cards aloud only with the presenter's permission; otherwise they take them home.

\begin{safety}[Showcase safety]
Only nicknames on screen and in file names. No uploading or sharing projects online. Photos or
videos only with every family's permission, and never with a screen showing personal details.
Remind students where online communities have age limits (often 13+) before they share code later.
\end{safety}

\begin{differentiation}
\supportpath Present together with a partner or with you as the \enquote{interviewer} asking the
prompt-card questions. Demo one feature only.
\extensionpath Explain an alternative design and its trade-off, or show a test that caught a bug.
\end{differentiation}
\end{document}
